\clearpage
\phantomsection
\label{backmatter:references}
\renewcommand{\refname}{References}
\begin{thebibliography}{127}
\bibitem{bib:source-1} R.A. Adams and J.J.F. Fournier. Sobolev spaces, volume 140 of Pure and Applied Mathematics (Amsterdam). Elsevier/Academic Press, Amsterdam, second edition, 2003.
\bibitem{bib:source-2} S. Agmon, S. Douglis, and L. Nirenberg. Estimates near the boundary for solutions of elliptic partial differential equations satisfying general boundary conditions I. Comm. Pure Appl. Math., 12:623–727, 1959.
\bibitem{bib:source-3} S. Agmon, S. Douglis, and L. Nirenberg. Estimates near the boundary for solutions of elliptic partial differential equations satisfying general boundary conditions II. Comm. Pure Appl. Math., 17:35–92, 1964.
\bibitem{bib:source-4} T. Alazard. A minicourse on the low Mach number limit. Discrete Contin. Dyn. Syst. Ser. S, 1(3):365–404, 2008.
\bibitem{bib:source-5} L. Ambrosio. Transport equation and Cauchy problem for BV vector fields. Invent. Math., 158(2):227–260, 2004.
\bibitem{bib:source-6} C. Amrouche, C. Bernardi, M. Dauge, and V. Girault. Vector potentials in three-dimensional non-smooth domains. Math. Methods Appl. Sci., 21(9):823–864, 1998.
\bibitem{bib:source-7} C. Amrouche and V. Girault. On the existence and regularity of the solution of Stokes problem in arbitrary dimension. Proc. Japan Acad. Ser. A Math. Sci., 67(5):171–175, 1991.
\bibitem{bib:source-8} M.E. Andersson. On the vector valued Hausdorff-Young inequality. Ark. Math., 36:1–30, 1998.
\bibitem{bib:source-9} P. Angot. Analysis of singular perturbation on the Brinkman problem for fictitious domain models of viscous flow. Math. Methods Appl. Sci., 22(15):1395–1412, 1999.
\bibitem{bib:source-10} P. Angot, C.H. Bruneau, and P. Fabrie. A penalization method to take into account obstacles in incompressible viscous flows. Numer. Math., 81:497–520, 1999.
\bibitem{bib:source-11} P. Angot and J.-P. Caltagirone. Natural convection through periodic porous media. In Proceeding 9th Int. Heat Transfer Conf., Jerusalem, Hemisphere,, volume 5, pages 219–224, 1990.
\bibitem{bib:source-12} S.N. Antontsev, A.V. Kazhikhov, and V.N. Monakhov. Boundary value problems in mechanics of nonhomogeneous fluids, volume 22 of Studies in Mathematics and Its Applications. North-Holland, Amsterdam, 1990.
\bibitem{bib:source-13} E. Arquis and J.-P. Caltagirone. Sur les conditions hydrodynamiques au voisinage d’une interface milieu fluide-milieu poreux: application à la convection naturelle. C. R. Acad. Sci Paris Série II, 299(1):1–4, 1984.
\bibitem{bib:source-14} T. Aubin. Un théorème de compacité. C. R. Acad. Sci. Paris, 256:5042–5044, 1963.
\bibitem{bib:source-15} R. Bellman. The stability of solutions of linear differential equations. Duke Math. J., 10:643–647, 1943.
\bibitem{bib:source-16} I. Bihari. A generalization of a lemma of Bellman and its application to uniqueness problems of differential equations. Acta Math. Acad. Sci. Hungar., 7:81–94, 1956.
\bibitem{bib:source-17} A. Blouza and H. Le Dret. An up-to-the-boundary version of Friedrichs’s lemma and applications to the linear Koiter shell model. SIAM J. Math. Anal., 33(4):877–895 (electronic), 2001.
\bibitem{bib:source-18} M.E. Bogovskiı̆. Solution of the first boundary value problem for the equation of continuity of an incompressible medium. Soviet. Math. Doklady, 20:1094–1098, 1979.
\bibitem{bib:source-19} M.E. Bogovskiı̆. Solutions of some problems of vector analysis, associated with the operators div and grad. In Theory of cubature formulas and the application of functional analysis to problems of mathematical physics, volume 1980 of Trudy Sem. S. L. Soboleva, No. 1, pages 5–40, 149. Akad. Nauk SSSR Sibirsk. Otdel. Inst. Mat., Novosibirsk, 1980.
\bibitem{bib:source-20} C. Bost, G.-H. Cottet, and E. Maitre. Convergence analysis of a penalization method for the three-dimensional motion of a rigid body in an incompressible viscous fluid. SIAM J. Numer. Anal., 48(4):1313–1337, 2010.
\bibitem{bib:source-21} J. Bourgain and H. Brezis. On the equation div Y = f and application to control of phases. J. Amer. Math. Soc., 16(2):393–426 (electronic), 2003.
\bibitem{bib:source-22} F. Boyer. Trace theorems and spatial continuity properties for the solutions of the transport equation. Diff. Integ. Eq., 18(8):891–934, 2005.
\bibitem{bib:source-23} F. Boyer. Analysis of the upwind finite volume method for general initial- and boundary-value transport problems. IMA J. Numer. Anal., 2011.
\bibitem{bib:source-24} F. Boyer and P. Fabrie. Éléments d’analyse pour l’étude de quelques modèles d’écoulements de fluides visqueux incompressibles, volume 52 of Mathématiques \& Applications (Berlin) [Mathematics \& Applications]. Springer-Verlag, Berlin, 2006.
\bibitem{bib:source-25} F. Boyer and P. Fabrie. Outflow boundary conditions for the incompressible non-homogeneous Navier-Stokes equations. Discrete Contin. Dyn. Syst. Ser. B, 7(2):219–250 (electronic), 2007.
\bibitem{bib:source-26} S.C. Brenner and L.R. Scott. The mathematical theory of finite element methods, volume 15 of Texts in Applied Mathematics. Springer, New York, third edition, 2008.
\bibitem{bib:source-27} H. Brezis. Functional analysis, Sobolev spaces and partial differential equations. Universitext. Springer, New York, 2011.
\bibitem{bib:source-28} C.-H. Bruneau. Numerical simulation and analysis of the transition to turbulence. In Paul Kutler, Jolen Flores, and Jean-Jacques Chattot, editors, Fifteenth International Conference on Numerical Methods in Fluid Dynamics, volume 490 of Lecture Notes in Physics, pages 1–13. Springer Berlin / Heidelberg, 1997.
\bibitem{bib:source-29} C.-H. Bruneau and P. Fabrie. Effective downstream boundary conditions for incompressible Navier-Stokes equation. Int. J. Num. Meth. Fluids, 19:693–705, 1994.
\bibitem{bib:source-30} C.-H. Bruneau and P. Fabrie. New efficient boundary conditions for incompressible Navier-Stokes equation: A well-posedness result. $M^2$AN, 30(7):815–840, 1996.
\bibitem{bib:source-31} C.-H. Bruneau and I. Mortazavi. Contrôle passif d’écoulements incompressibles autour d’obstacles à l’aide de milieux poreux. C. R. Acad. Sci. Paris Série II, 329:517–521, 2001.
\bibitem{bib:source-32} V. Bruneau and V. Carbou. Asymptotic expansion for large limit coupling. Asymp. Anal., 29(2):91–113, 2002.
\bibitem{bib:source-33} V.I. Burenkov. Sobolev spaces on domains, volume 137 of Teubner-Texte zur Mathematik [Teubner Texts in Mathematics]. Teubner Verlagsgesellschaft mbH, Stuttgart, 1998.
\bibitem{bib:source-34} J.-P. Caltagirone. Sur l’interaction fluide-milieux poreux : application au calcul des efforts exercés sur un obstacle par un fluide visqueux. C. R. Acad. Sci. Paris Série II, 318:571–577, 1994.
\bibitem{bib:source-35} G. Carbou. Penalization method for viscous incompressible flow around a porous thin layer. Nonlinear Anal. Real World Appl., 5(5):815–855, 2004.
\bibitem{bib:source-36} G. Carbou. Brinkmann model and double penalization method for the flow around a porous thin layer. J. Math. Fluid Mech., 10(1):126–158, 2008.
\bibitem{bib:source-37} G. Carbou and P. Fabrie. Boundary layers for a penalisation method for incompressible flows. Adv. Diff. Eq., 8(12):1453–1480, 2003.
\bibitem{bib:source-38} G. Carbou, P. Fabrie, and O. Guès. Couches limites dans un modèle de ferromagnétisme. Comm. Partial Diff. Eq., 27(7–8):1467–1495, 2002.
\bibitem{bib:source-39} L. Cattabriga. Su un problema al contorno relativo al sistema di equazioni di Stokes. Rend. Sem. Padova, 31:308–340, 1961.
\bibitem{bib:source-40} S. Chandrasekhar. Hydrodynamic and hydromagnetic stability. Dover, New York, 1961.
\bibitem{bib:source-41} J.-Y. Chemin. Perfect incompressible fluids, volume 14 of Oxford Lecture Series in Mathematics and Its Applications. Clarendon Press, Oxford University Press, New York, 1998.
\bibitem{bib:source-42} H. Choe and H. Kim. Strong solutions of the Navier-Stokes equations for nonhomogeneous incompressible fluids. Comm. Partial Differential Equations, 28(5-6):1183–1201, 2003.
\bibitem{bib:source-43} P. Chossat and G. Iooss. The Couette-Taylor problem. Number 102 in Applied Mathematical Sciences. Springer Verlag, New York, 1994.
\bibitem{bib:source-44} P. Constantin and C. Foias. Navier-Stokes equations. Chicago Lectures in Mathematics. University of Chicago Press, 1988.
\bibitem{bib:source-45} R. Danchin. Density-dependent incompressible viscous fluids in critical spaces. Proc. Roy. Soc. Edinburgh Sect. A, 133(6):1311–1334, 2003.
\bibitem{bib:source-46} R. Danchin. Low Mach number limit for viscous compressible flows. $M^2$AN Math. Model. Numer. Anal., 39(3):459–475, 2005.
\bibitem{bib:source-47} R. Dautray and J.-L. Lions. Mathematical analysis and numerical methods for science and technology. Vol. 3. Springer-Verlag, Berlin, 1990. Spectral theory and applications.
\bibitem{bib:source-48} R. Dautray and J.-L. Lions. Mathematical analysis and numerical methods for science and technology. Vol. 5. Springer-Verlag, Berlin, 1992. Evolution problems. I.
\bibitem{bib:source-49} G. de Rham. Differentiable manifolds, volume 266 of Grundlehren der Mathematischen Wissenschaften [Fundamental Principles of Mathematical Sciences]. Springer-Verlag, Berlin, 1984. Forms, currents, harmonic forms,.
\bibitem{bib:source-50} B. Desjardins. Global existence results for the incompressible density-dependent Navier-Stokes equations in the whole space. Diff. Integral Eq., 10:587–598, 1997.
\bibitem{bib:source-51} B. Desjardins. Linear transport equations with initial values in Sobolev spaces and application to the Navier-Stokes equations. Diff. Integral Eq., 10:577–586, 1997.
\bibitem{bib:source-52} R.J. Di Perna and P.-L. Lions. Ordinary differential equations, transport theory and Sobolev spaces. Invent. Math., 98:511–547, 1989.
\bibitem{bib:source-53} P.G. Drazin and W.H. Reid. Hydrodynamic stability. Cambridge Monographs on Mechanics and Applied Mathematics. Cambridge University Press, Cambridge, New York, 1982.
\bibitem{bib:source-54} G. Duvaut and J.-L. Lions. Inequalities in mechanics and physics. Springer-Verlag, Berlin, 1976.
\bibitem{bib:source-55} A. Ern and J.-L. Guermond. Theory and practice of finite elements, volume 159 of Applied Mathematical Sciences. Springer-Verlag, New York, 2004.
\bibitem{bib:source-56} L.C. Evans. Partial differential equations, volume 19 of Graduate Studies in Mathematics. American Mathematical Society, Providence, RI, second edition, 2010.
\bibitem{bib:source-57} E. Feireisl. Dynamics of viscous compressible fluids, volume 26 of Oxford Lecture Series in Mathematics and its Applications. Oxford University Press, Oxford, 2004.
\bibitem{bib:source-58} E. Feireisl. Mathematical models of incompressible fluids as singular limits of complete fluid systems. Milan J. Math., 78(2):523–560, 2010.
\bibitem{bib:source-59} C. Foias, O. Manley, R. Rosa, and R. Temam. Navier-Stokes equations and turbulence. Encyclopedia of Mathematics and Its Applications. Cambridge University Press, Cambridge, 2001.
\bibitem{bib:source-60} C. Foias and R. Temam. Remarques sur les équations de Navier-Stokes et les phénomènes successifs de bifurcation. Annal. Sc. Norm. Sup. Pisa, Serie IV, 5:29–63, 1978.
\bibitem{bib:source-61} K.O. Friedrichs. The identity of weak and strong extensions of differential operators. Trans. Amer. Math. Soc., 55:132–151, 1944.
\bibitem{bib:source-62} G.P. Galdi. An introduction to the Navier-Stokes initial-boundary value problem. In Fundamental directions in mathematical fluid mechanics, Adv. Math. Fluid Mech., pages 1–70. Birkhäuser, Basel, 2000.
\bibitem{bib:source-63} G.P. Galdi. An introduction to the mathematical theory of the Navier-Stokes equations. Springer Monographs in Mathematics. Springer, New York, second edition, 2011. Steady-state problems.
\bibitem{bib:source-64} J. Garcia-Cuerva, K.S. Kazarian, V.I. Kolyada, and J.L. Torrea. Vector-valued Hausdorff-Young inequality and applications. Russ. Math. Surv., 53(3):435–513, 1998.
\bibitem{bib:source-65} V. Girault and P.-A. Raviart. Finite element methods for Navier-Stokes equations. Springer-Verlag, 1986.
\bibitem{bib:source-66} O. Gonzalez and A. M. Stuart. A first course in continuum mechanics. Cambridge Texts in Applied Mathematics. Cambridge University Press, Cambridge, 2008.
\bibitem{bib:source-67} P. Grisvard. Elliptic problems in nonsmooth domains. Monographs and Studies in Mathematics, 24. Pitman Advanced Publishing Program. Pitman, Boston-London-Melbourne XIV, 410 p., 1985.
\bibitem{bib:source-68} J. Heinonen. Lectures on Lipschitz analysis, volume 100 of Report. University of Jyväskylä Department of Mathematics and Statistics. University of Jyväskylä, Jyväskylä, 2005.
\bibitem{bib:source-69} E. Hewitt and K. Stromberg. Real and Abstact Analysis. Springer Verlag, New York,, 1965.
\bibitem{bib:source-70} J.G. Heywood. On the impossibility, in some cases, of the Leray-Hopf condition for energy estimates. J. Math. Fluid Mech., 13(3):449–457, 2011.
\bibitem{bib:source-71} E. Hopf. Ein allgemeiner Endlichkeitssatz der Hydrodynamik. Math. Ann., 117:764–775, 1941.
\bibitem{bib:source-72} E. Hopf. Über die Anfangswertaufgabe für die hydrodynamischen Grundgleichungen. Math. Nachr., 4:213–231, 1951.
\bibitem{bib:source-73} R.B. Kellogg and J.E. Osborn. A regularity result for the Stokes problem in a convex polygon. J. Funct. Anal., 21:397–431, 1976.
\bibitem{bib:source-74} N. Kevlahan and J.-M. Ghidaglia. Computation of turbulent flow past an array of cylinders using a spectral method with Brinkman penalization. Eur. J. Mech., B, Fluids, 20(3):333–350, 2001.
\bibitem{bib:source-75} K. Khadra, S. Parneix, P. Angot, and J.-P. Caltagirone. Fictitious domain approach for numerical modeling of Navier-Stokes equation. Int. J. Numer. Meth. in Fluids, 34(8):651–684, 2000.
\bibitem{bib:source-76} O. Ladyzhenskaya. The Mathematical Theory of Viscous Incompressible Flow. Gordon and Breach, New York, 1963.
\bibitem{bib:source-77} O. Ladyzhenskaya and V.A. Solonnikov. Unique solvability of an initial and boundary value problem for viscous incompressible non-homogeneous fluids. J. ZSoviet. Math., 9:697–749, 1978.
\bibitem{bib:source-78} C. Le Bris and P.-L. Lions. Renormalized solutions of some transport equations with partially $W^{1,1}$ velocities and applications. Ann. Mat. Pura Appl. (4), 183(1):97–130, 2004.
\bibitem{bib:source-79} J. Leray. Essai sur les mouvements plans d’un liquide visqueux emplissant l’espace. Acta Math., 63:193–248, 1934.
\bibitem{bib:source-80} J. Leray. Essai sur les mouvements plans d’un liquide visqueux que limitent des parois. J. Math. Pures Appl., 13:331–418, 1934.
\bibitem{bib:source-81} R. Lewandowski and G. Pichot. Numerical simulation of water flow around a rigid fishing net. Comput. Methods Appl. Mech. Engrg., 196(45-48):4737–4754, 2007.
\bibitem{bib:source-82} G.M. Lieberman. Regularized distance and its applications. Pacific J. Math., 117(2):329–352, 1985.
\bibitem{bib:source-83} T. Lining and Y. Hongjun. Classical solutions to Navier-Stokes equations for nonhomogeneous incompressible fluids with non-negative densities. J. Math. Anal. Appl., 362(2):476–504, 2010.
\bibitem{bib:source-84} J.-L. Lions. Quelques méthodes de résolution des problèmes aux limites non linéaires. Dunod, Paris, 1969.
\bibitem{bib:source-85} J.-L. Lions and E. Magenes. Problèmes aux limites non homogènes et applications. Vol. 1. Travaux et Recherches Mathématiques, No. 17. Dunod, Paris, 1968.
\bibitem{bib:source-86} P.-L. Lions. Mathematical topics in fluid mechanics, Tome 1 : Incompressible models, volume 3 of Oxford Lecture Series in Mathematics and Applications. Oxford University Press, New York, 1996.
\bibitem{bib:source-87} P.-L. Lions. Mathematical topics in fluid mechanics, Tome 2 : Compressible models, volume 10 of Oxford Lecture Series in Mathematics and Applications. Oxford University Press, New York, 1998.
\bibitem{bib:source-88} A.J. Majda and A.L. Bertozzi. Vorticity and incompressible flow, volume 27 of Cambridge Texts in Applied Mathematics. Cambridge University Press, Cambridge, UK, 2002.
\bibitem{bib:source-89} V.G. Maz’ya and S.V. Poborchi. Differentiable functions on bad domains. World Scientific, River Edge, NJ, 1997.
\bibitem{bib:source-90} R.E. Meyer. Introduction to mathematical fluid dynamics. Dover, New York, 1982. Corrected reprint of the 1971 original.
\bibitem{bib:source-91} J. Nečas. Sur les normes equivalentes dans $W_p^k$ et sur la coercivite des formes formellement positives. Sem. math. sup. Universite Montreal, pages 102–128, 1966.
\bibitem{bib:source-92} J. Nečas. Direct Methods in the Theory of Elliptic Equations. Springer, New York, 2012.
\bibitem{bib:source-93} A. Novotný and I. Straškraba. Introduction to the mathematical theory of compressible flow, volume 27 of Oxford Lecture Series in Mathematics and its Applications. Oxford University Press, Oxford, 2004.
\bibitem{bib:source-94} J.P. O’Connell and J.M. Haile. Thermodynamics: Fundamentals for Applications (Cambridge Series in Chemical Engineering). Cambridge University Press, Cambridge, UK,, 2005.
\bibitem{bib:source-95} A. Pazy. Semigroups of linear operators and applications to partial differential equations, volume 44 of Applied Mathematical Sciences. Springer-Verlag, New York, 1983.
\bibitem{bib:source-96} J. Peetre. Sur la transformation de Fourier des fonctions à valeurs vectorielles. Rend. Sem. Math. Univ. Padova, 42:15–26, 1969.
\bibitem{bib:source-97} C. Peskin. Numerical analysis of blood flow in the heart. J. Comput. Phys., 25:220–252, 1977.
\bibitem{bib:source-98} C. Peskin. The fluid dynamics of heart valves: Experimental, theoretical and comptutational method. Annu. Rev. Fluid Mech, 14:235–259, 1982.
\bibitem{bib:source-99} W. Rudin. Principles of mathematical analysis. International Series in Pure and Applied Mathematics. McGraw-Hill, New York, third edition, 1976.
\bibitem{bib:source-100} W. Rudin. Real and complex analysis. McGraw-Hill, New York, third edition, 1987.
\bibitem{bib:source-101} W. Rudin. Functional analysis. International Series in Pure and Applied Mathematics. McGraw-Hill, New York, second edition, 1991.
\bibitem{bib:source-102} I.L. Ryhming. Dynamique des fluides. Presses Polytechniques et Universitaires Romandes, Lausanne, second edition, 2004. Un cours de base du deuxième cycle universitaire.
\bibitem{bib:source-103} E.M. Saiki and S. Biringen. Numerical simulation of a cylinder in uniform flow:application of a virtual boundary method. J. Comput. Phys., 123:450–465, 1996.
\bibitem{bib:source-104} L. Schwartz. Analyse. Deuxième partie: Topologie générale et analyse fonctionnelle. Collection Enseignement des Sciences, No. 11. Hermann, Paris, 1970.
\bibitem{bib:source-105} L. Schwartz. Analyse. I Théorie des ensembles et topologie. [Set theory and topology], With the collaboration of K. Zizi, volume 42 of Collection Enseignement des Sciences [Collection: The Teaching of Science]. Hermann, Paris, 1991.
\bibitem{bib:source-106} J. Serrin. On the interior regularity of weak solutions of the Navier-Stokes equations. Arch. Rational Mech. Anal., 9:187–195, 1962.
\bibitem{bib:source-107} J. Serrin. The initial value problem for the Navier-Stokes equations. In Nonlinear Problems (Proc. Sympos., Madison, Wis., 1962), pages 69–98. Univ. of Wisconsin Press, Madison, 1963.
\bibitem{bib:source-108} J. Simon. Caractérisation d’espaces fonctionnels. Bolletino Un. Math. Ital. B, 5(2):687–714, 1978.
\bibitem{bib:source-109} J. Simon. Compact sets in the space $L^p(0,T;B)$. Ann. Mat. Pura appl., 146(4):65–96, 1987.
\bibitem{bib:source-110} J. Simon. Nonhomogeneous viscous incompressible fluids: Existence of velocity, density and pressure. SIAM J. Math. Anal., 21(5):1093–1117, 1990.
\bibitem{bib:source-111} J. Simon. Démonstration constructive d’un théorème de G. de Rham. C. R. Acad. Sci. Paris, Sér. I, 316:1167–1172, 1993.
\bibitem{bib:source-112} J. Simon. Representation of distributions with explicit antiderivatives up to the boundary. In Progress in partial differential equations: The Metz surveys, 2 (1992), volume 296 of Pitman Res. Notes Math. Ser., pages 201–215. Longman Sci. Tech., Harlow, 1993.
\bibitem{bib:source-113} J. Simon. Distributions et espaces de Lebesgue à valeurs dans un espace vectoriel topologique localement convexe, 2009. \url{http://math.unice.fr/~jsimon/pubs/Simon-E12.pdf}.
\bibitem{bib:source-114} D.R. Smart. Fixed point theorems, volume 66 of Cambridge tracts in mathematics. Cambridge University Press, 1974.
\bibitem{bib:source-115} H. Sohr. Zur Regularitätstheorie der instationären Gleichungen von Navier-Stokes. Math. Z., 184(3):359–375, 1983.
\bibitem{bib:source-116} H. Sohr and W. von Wahl. On the singular set and the uniqueness of weak solutions of the Navier-Stokes equations. Manuscripta Math., 49(1):27–59, 1984.
\bibitem{bib:source-117} E. M. Stein. Singular integrals and differentiability properties of functions. Princeton Mathematical Series, No. 30. Princeton University Press, Princeton, NJ, 1970.
\bibitem{bib:source-118} L. Tartar. Topics in nonlinear analysis, volume 13 of Publications Mathématiques d’Orsay 78. Université de Paris-Sud Département de Mathématique, Orsay, 1978.
\bibitem{bib:source-119} L. Tartar. An introduction to Navier-Stokes equation and oceanography, volume 1 of Lecture Notes of the Unione Matematica Italiana. Springer-Verlag, Berlin, 2006.
\bibitem{bib:source-120} R. Temam. Navier-Stokes equations and nonlinear functional analysis, volume 66 of CBMS-NSF Regional Conference Series in Applied Mathematics. Society for Industrial and Applied Mathematics (SIAM), Philadelphia, second edition, 1995.
\bibitem{bib:source-121} R. Temam. Infinite Dimensional Dynamical Systems in Mechanics and Physics. Second edition, volume 68 of Applied Mathematical Sciences. Springer Verlag, New York, 1997.
\bibitem{bib:source-122} R. Temam. Navier-Stokes equations. AMS Chelsea, Providence, RI, 2001. Theory and numerical analysis, Reprint of the 1984 edition.
\bibitem{bib:source-123} R. Temam and A. Miranville. Mathematical modeling in continuum mechanics. Cambridge University Press, Cambridge, UK, second edition, 2005.
\bibitem{bib:source-124} X. M. Wang. A remark on the characterization of the gradient of a distribution. Appl. Anal., 51(1-4):35–40, 1993.
\bibitem{bib:source-125} C. Wassgren, C. M. Krousgrill, and P. Carmody. The continuum assumption. \url{https://engineering.purdue.edu/~wassgren/applet/java/continuum/index.html}.
\bibitem{bib:source-126} R. Kh. Zeytounian. Theory and applications of nonviscous fluid flows. Springer-Verlag, Berlin, 2002.
\bibitem{bib:source-127} R. Kh. Zeytounian. Theory and applications of viscous fluid flows. Springer-Verlag, Berlin, 2004.
\end{thebibliography}
